\documentclass[10pt,a4paper]{amsart}
\usepackage[margin=19mm]{geometry}
\usepackage{amsmath,amssymb,mathtools}
\usepackage[scr=rsfs,cal=euler]{mathalfa}
\usepackage{xfrac}
\usepackage[unicode,hidelinks,bookmarks=false]{hyperref}
\newcommand{\ie}{i.e.\ }
\newcommand{\cf}{cf.\ }
\newcommand{\resp}{respectively\ }
\newcommand{\Subsection}{\subsection}
\providecommand{\nobreakdash}{\nobreak\hbox{-}}
\setlength{\emergencystretch}{2em}
\multlinegap0pt
\usepackage[shortlabels]{enumitem}
\usepackage{amssymb}
\newtheorem{theorem}{Theorem}
\newtheorem{proposition}{Proposition}
\newtheorem{lemma}{Lemma}
\newtheorem{corollary}{Corollary}
\title{A uniform decomposition theorem for invariant differential operators on imprimitive complex reflection groups $G(r,p,n)$}
\author{Jean Kabor\'e,, Ibrahim Nonkan\'e}
\date{Source: arXiv:2608.01474v1 (CC BY 4.0)}
\begin{document}
\maketitle

\noindent\emph{Source and licence.} A uniform decomposition theorem for invariant differential operators on imprimitive complex reflection groups $G(r,p,n)$. Version arXiv:2608.01474v1 (CC BY 4.0), distributed by arXiv under the Creative Commons Attribution 4.0 International licence, http://creativecommons.org/licenses/by/4.0/, as recorded in the arXiv licence metadata for this exact version and shown on the abstract page https://arxiv.org/abs/2608.01474v1. Full licence notice: \texttt{licenses/arxiv-2608.01474-CC-BY-4.0.txt}.\par\smallskip

\noindent\textit{Editorial scope.} Counted unit: the article's corollary cor:twisted-invariants (source0.tex lines 1229--1237). The article's complete original proof of it is delivered with it (source0.tex lines 1239--1296, one \texttt{proof} environment of 58 lines). No mathematics is written, repaired or re-derived: the statement and the proof are the article's own lines, in the article's own notation, and no retained line is substituted or re-flowed.  Retained uncounted as the source-local context the proof needs: lines 791--806; lines 852--872; lines 1053--1061; lines 1110--1117; lines 1165--1173.  Nothing was omitted from any retained span, so the package carries no omission record.  No retained line contains a citation, so the wrapper carries no bibliography and no bibliography entry.  No retained line contains a figure environment, an image-inclusion command or drawing code, so no image is delivered and the package declares no asset dependency.  Editorial formatting only, in addition: an AMS \texttt{amsart} preamble that restates the article's own declarations and macros used by the retained lines, namely the environments \texttt{theorem}, \texttt{proposition}, \texttt{lemma}, \texttt{corollary} on separate counters, together with the standard packages those lines need.  The retained environments print in this document as Theorem 1, Proposition 1, Lemma 1, Corollary 1 in order of appearance, whereas the article prints the same items with the numbering of its own full document.  The complete native source of the article as distributed in the arXiv e-print for this exact version is bundled unchanged as \texttt{sources/206632/source0.tex}.
\begin{theorem}\label{thm:main-general}
Fix $\lambda \in \mathcal{P}_{r,n}/\sim$ and $0 \le l < e(\lambda)$,
and enumerate $\mathrm{STab}(\lambda)_d = \{S_1,\dots,S_{f^l_\lambda}\}$.
For every $i \in \{1,\dots,f^l_\lambda\}$, writing $e^{\lambda,l}_i$
for the corresponding primitive idempotent of $\mathbb{C}[W]$
(as described above):
\begin{enumerate}[label=(\roman*)]
 \item $e^{\lambda,l}_i\,\widetilde{\mathcal{O}}_X$ is a nontrivial
 $\widetilde{\mathcal{D}}_Y$-submodule of $\widetilde{\mathcal{O}}_X$;
 \item the $\widetilde{\mathcal{D}}_Y$-module
 $e^{\lambda,l}_i\,\widetilde{\mathcal{O}}_X$ is simple;
 \item there exists a higher Specht polynomial $F^l_{S_i,T}$, for some
 $T \in \mathrm{STab}(\lambda)_{d\,b(\lambda)}$, such that
 $e^{\lambda,l}_i\,\widetilde{\mathcal{O}}_X = \widetilde{\mathcal{D}}_Y \, F^l_{S_i,T}$.
\end{enumerate}
\end{theorem}

\begin{proposition}\label{prop:decomposition-general}
With the enumeration $\mathrm{STab}(\lambda)_d = \{S_1,\dots,S_{f^l_\lambda}\}$
of Theorem~\ref{thm:main-general}, write
$F^l_{\lambda,i} := F^l_{S_i,T}$ for the generator of
$e^{\lambda,l}_i\,\widetilde{\mathcal O}_X$ given by
Theorem~\ref{thm:main-general}(iii). Then:
\begin{enumerate}[label=(\roman*)]
\item
\[
 \widetilde{\mathcal{O}}_X
 = \bigoplus_{\lambda \in \mathcal{P}_{r,n}/\sim} \bigoplus_{l=0}^{e(\lambda)-1}
 \bigoplus_{i=1}^{f^l_\lambda} \widetilde{\mathcal{D}}_Y F^l_{\lambda,i} ;
\]
\item fixing, for each $(\lambda,l)$, the representative $i=1$,
\[
 \widetilde{\mathcal{O}}_X
 = \bigoplus_{\lambda \in \mathcal{P}_{r,n}/\sim} \bigoplus_{l=0}^{e(\lambda)-1}
 f^l_\lambda \, \widetilde{\mathcal{D}}_Y F^l_{\lambda,1}.
\]
\end{enumerate}
\end{proposition}

\begin{lemma}[Descent isomorphism]\label{lem:descent-iso}
For every finite-dimensional $\mathbb{C}$-vector space $V$ with a
$W$-action, the multiplication map
\[
  \mu_V : L \otimes_K \big(L \otimes_{\mathbb{C}} V\big)^{W} \longrightarrow L \otimes_{\mathbb{C}} V,
  \qquad a \otimes m \longmapsto am
\]
is an isomorphism of $L$-vector spaces.
\end{lemma}

\begin{lemma}[The extended partials span $\operatorname{Der}_{\mathbb{C}}(L)$]\label{lem:der-span}
The $L$-vector space $\operatorname{Der}_{\mathbb{C}}(L)$ of
$\mathbb{C}$-linear derivations of $L$ is free of rank $n$, with basis
given by the unique extensions
$\widetilde{\partial/\partial y_1},\dots,\widetilde{\partial/\partial y_n}$
of the derivations $\partial/\partial y_1,\dots,\partial/\partial y_n$
of $K$ furnished by (F1).
\end{lemma}

\begin{lemma}[Rigidity of flat endomorphisms]\label{lem:rigidity}
Let $V$ be a finite-dimensional $\mathbb{C}$-vector space and equip
$L \otimes_{\mathbb{C}} V$ with the connection
$\widetilde{D} \otimes \operatorname{id}$ for $\widetilde{D}$ ranging
over derivations of $L$ extending those of $K$. Then any $L$-linear
endomorphism of $L \otimes_{\mathbb{C}} V$ commuting with
$\widetilde{D} \otimes \operatorname{id}$ for all such $\widetilde{D}$
lies in $\operatorname{id}_L \otimes \operatorname{End}_{\mathbb{C}}(V)$.
\end{lemma}

\begin{corollary}[Twisted-invariants description of the simple summands]\label{cor:twisted-invariants}
For every $\lambda \in \mathcal{P}_{r,n}/\sim$ and $0 \le l < e(\lambda)$,
\[
  N_\lambda := \widetilde{\mathcal{D}}_Y F^l_{\lambda,1}
  \;\cong\; \nabla\big(V^l(\lambda)\big)
  \;=\; \big(\widetilde{\mathcal{O}}_X \otimes_{\mathbb{C}} V^l(\lambda)\big)^{W}
\]
as $\widetilde{\mathcal{D}}_Y$-modules.
\end{corollary}

\begin{proof}
We first show $\nabla(V^l(\lambda))$ is simple. Let $M \subset \nabla(V^l(\lambda))$
be a nonzero $\widetilde{\mathcal{D}}_Y$-submodule and set
$\overline M := L \otimes_{\widetilde{\mathcal{O}}_Y} M$, a nonzero
$L$-subspace of
$L \otimes_{\widetilde{\mathcal{O}}_Y} \nabla(V^l(\lambda)) \cong L \otimes_{\mathbb{C}} V^l(\lambda)$
(Lemma~\ref{lem:descent-iso}) that is stable under
$\widetilde D \otimes \operatorname{id}$ for every derivation
$\widetilde D$ of $L$ extending one of $K$, since $M$ is a
$\widetilde{\mathcal{D}}_Y$-submodule. Write $V := V^l(\lambda)$ and
fix a $\mathbb{C}$-basis $v_1,\dots,v_f$ of $V$. Pick
$0 \ne w = \sum_i a_i \otimes v_i \in \overline M$ with a minimal
number $r$ of nonzero coefficients among all nonzero elements of
$\overline M$; relabel so $a_1,\dots,a_r \ne 0$, and scale (using that
$\overline M$ is an $L$-subspace) so $a_1 = 1$. For any extended
derivation $\widetilde D$, stability gives
$\widetilde D(w) = \sum_i \widetilde D(a_i) \otimes v_i \in \overline M$;
since $\widetilde D(a_1) = \widetilde D(1) = 0$, this element has at
most $r-1$ nonzero coefficients, so by minimality of $r$ it must be
$0$ -- exactly the minimal-support argument of
Lemma~\ref{lem:descent-iso}'s injectivity proof, now applied to a
subspace rather than a relation. Hence $\widetilde D(a_i) = 0$ for
every $i$ and every such $\widetilde D$; by Lemma~\ref{lem:der-span}
these span $\operatorname{Der}_{\mathbb{C}}(L)$ over $L$, so
$D(a_i) = 0$ for \emph{every} $\mathbb{C}$-derivation $D$ of $L$, and
as in the proof of Lemma~\ref{lem:rigidity} this forces $a_i \in \mathbb{C}$
for every $i$. Thus $w = 1 \otimes u$ for the nonzero vector
$u := \sum_i a_i v_i \in V$.

Let $U := \{u \in V : 1 \otimes u \in \overline M\}$, a
$\mathbb{C}$-subspace with $U \ne 0$ by the above; clearly
$L \otimes_{\mathbb{C}} U \subset \overline M$. If this inclusion were
strict, the same argument applied to a nonzero element of
$\overline M$ chosen outside $L \otimes_{\mathbb{C}} U$ (with minimal
support relative to a basis of $V$ extending one of $U$) would
produce a further nonzero constant vector in $U$ not already
accounted for, contradicting finite-dimensionality after finitely
many repetitions; hence $\overline M = L \otimes_{\mathbb{C}} U$.
Since $V$ is a simple $\mathbb{C}[W]$-module and $M$, hence
$\overline M$, hence $U$, is $W$-stable (as $M$ is a
$\widetilde{\mathcal{D}}_Y$-submodule of $\nabla(V)$, and the
identification with $L\otimes_{\mathbb C}V$ is $W$-equivariant), $U$
is a nonzero $W$-submodule of $V$, so $U = V$. Thus
$\overline M = L \otimes_{\mathbb{C}} V$, i.e.\ $M$ has full rank
$\dim_{\mathbb{C}} V$ inside $\nabla(V)$, which itself has rank
$\dim_{\mathbb{C}} V$ (Lemma~\ref{lem:descent-iso}); a submodule of
full rank in a torsion-free module of the same rank, with both sides
finitely generated over the same ring, forces $M = \nabla(V)$.
Hence $\nabla(V^l(\lambda))$ is simple.

The $\widetilde{\mathcal{O}}_Y$-rank of
$\nabla(V^l(\lambda))$ equals $\dim_{\mathbb{C}} V^l(\lambda) = f^l_\lambda$,
matching the rank of $N_\lambda$ (Theorem~\ref{thm:main-general}(iii),
Proposition~\ref{prop:decomposition-general}), and both are simple
$\widetilde{\mathcal{D}}_Y$-modules whose associated $W$-representation
(via $\operatorname{Loc}$) is $V^l(\lambda)$. By the uniqueness in
Theorem~\ref{thm:main-general}(ii)--(iii), $N_\lambda \cong \nabla(V^l(\lambda))$.
\end{proof}

\end{document}
